% !TEX root = main.tex
\section{보안 요구사항}
본 논문의 보안 성질은 정확성(Correctness), 건전성(Sound-
ness), 프라이버시(Privacy) 세 가지다. 정상 실행이 의도한 결
과를 내는지, 잘못된 결과를 성공으로 받아들이게 할 수 있는지,
허용된 공개 정보 외의 비밀이 숨겨지는지를 각각 다룬다. 추적
·동결과 감사키 보호도 이 세 성질 안에서 분석한다. 다음 절의
정리는 명시한 계약 아래의 결과이며 실제 암호 방식·회로·계약
의 가정 충족 여부와 구분한다.

\subsection{정확성(Correctness)}
필요한 정상 조건이 갖춰졌을 때 정직한 실행은 기대 결과를 내
야 한다. 유효한 전이는 전체 실행 조건, 출력 신규성과 저장 용
량을 만족하는 안정된 상태에서 수락되어 규범의 상태 갱신을
수행한다. 발급이 포함된 인증된 이력에서 그 발급과 일치하는
정규 문서·제품 주장은 공개 DPP 검증을 통과한다.

유효한 시작점과 확정 \(H_{t}\), 외부에서 확보한 유효 감사키, 필
요한 유한 기록과 종료하는 인증 조회가 있으면 선행·후속 추적
은 정확한 그래프·경계·종단을 반환해야 한다. 고정 \(H_{t}\)의 추적
에 그 이후 원장 전체의 무경합을 요구하지는 않는다.

차단 승인이 있는 경우 정확한 전체 대상 \(T\) 를 처리한다.
\(T\) 가 비면 NO\_LIVE\_TARGETS다. \(T\) 가 비어 있지 않고 대상
소비·차단 해제 경합이 없으며, 필요한 동결·조회·제출 처리
와 공통 시점 확인이 설정된 유한 한도 안에 끝나면 COM-
PLETE\_AT\_CHECKPOINT를 반환해야 한다. 자료·키 부재, 소
비 경합과 확정 지연이 있는 상황까지 정상 성공을 보장하는 성
질은 아니다.

% Source PDF page 11, column 2

\subsection{건전성(Soundness)}
공격자가 실제 생성 값에 맞지 않는 전이를 수락시키거나 잘못
된 추적·차단 결과를 성공으로 보고하게 하기 어려워야 한다.
전이에서는 개방값·포함 관계·소유권·정수 회계·권한·정책의
위반, 반복 계상·소비, 현재 상태·기한의 위반을 배제한다. 감사
암호문은 실제 부모와 미래 소비 식별자에 결합되어야 하며, 미
등록 또는 발급 기록과 불일치하는 공개 DPP 주장이 유효로 판
정되어서는 안 된다.

감사에서는 누락 기록이나 조회·복호화 오류를 빈 경로나
미소비로 해석해 잘못된 성공을 반환해서는 안 된다. 동결
에서는 전체 대상을 보존하고 같은 확정 시점에서 완료를 확
인해야 한다. 부분 결과, 소비 경합과 불명확한 제출을 남긴
INCOMPLETE는 오류가 아니다. 이를 전체 완료로 바꾸어 보
고하는 것이 건전성 위반이다. 전체 \(T\) 가 공통 시점에 미소비·차
단이면, 그 차단이 연속 유지되는 동안 후속 소비도 수락되어서
는 안 된다.

이 성질은 증명에 의한 실제 값의 결합과 원장의 원자적 현재
상태 집행을 함께 필요로 한다. 신뢰된 감사자·상태 관리자와
승인된 관계에 대한 보장이며 외부 측정값의 진실성이나 정책
의 현실적 적합성을 교정하는 결과는 아니다. 영구 차단·기발급
주장 취소·물리적 회수도 포함하지 않는다.

\subsection{프라이버시(Privacy)}
숨기려는 정보는 비공개 개방값·소유 비밀·회계 수치·실제 부
모와 생성 참조에서 미래 소비 식별자로 가는 대응이다. 허용된
공개 의미가 같은 두 유효 이력의 전체 공개 관측 분포를 구별하
기 어려워야 한다. 공개 의미에는 종류·정책·적용 범위·입출력
형상·sibling 관계·시간· 호출자·기한·문서와 제품 주장의 의미,
공개 색인의 동일성 패턴이 포함된다. 비교 이력의 커미트먼트
나 암호문 바이트 자체가 같아야 한다는 뜻은 아니다.

데이터 공개키 \(PK_{B}\)의 기록 암호화와 임계 공개키 \(PK_{A}\)의
\(SK_{B}\) 보호는 프라이버시의 일부다. 공개 키와 \(W_{B}\), 설정 전
에 고정한 \(K−1\)개 이하 위원회 조각을 관측에 포함한다. 초
기 비밀 사본 관리와 외부 키 복구·전달은 III절의 신뢰 경계다.
\(W_{B}\)를 공개하는 것과 \(SK_{B}\)를 공개하는 것은 다르며, \(SK_{B}\)를
받은 주체의 동일 키 영역 복호화를 제한하지 않는다.

현재 증명하는 범위는 정직하게 생성된 이력을 읽는 수동 외
부 관찰자다. 참여자 비밀·개방값, 보호 관계의 감사·상태 공개
와 능동 거래·복호화 질의는 게임에서 제외한다. 승인된 정보
공개까지 허용하는 일반 요구에는 공개된 식별자의 추가 연결
과 참여자의 기존 지식을 반영한 별도 게임이 필요하다. 외부
DPP의 공개 내용이나 한 번 공개된 관계를 다시 숨기는 성질은
아니다. 두 키 암호화의 기밀성만으로 커미트먼트·소비 식별자·
증명의 공동 분포가 숨겨진다고 추정하지 않는다. 상세 게임과
환원은 보조 문서 zkdpp-v2-security-proofs.md에 정리
한다.
